% from main.tex
%%%%%%%%%%  scan idx 393  |  folio 380  %%%%%%%%%%
\begin{equation}\tag{1.13.2}
  x^{r+1}=0.
\end{equation}

\textbf{Remarque 1.14.} En utilisant la notion de «schéma relatif» (cf [3]), on peut encore donner un sens à 1.1 en supposant que \(S\) est un topos localement annelé. Si on suppose que l'objet final du topos est quasi-compact, la démonstration donnée dans le cas des schémas se transpose sans difficultés, et le théorème reste donc valable.

\section*{2. Calcul du \(K^{\bullet}\) d'un fibré projectif : cas des complexes parfaits}


Le présent paragraphe ne sera pas utilisé dans la suite du séminaire. La notation \(K^{\bullet}(X)\) désignera dans ce paragraphe le groupe de Grothendieck des complexes parfaits sur \(X\) (IV 2.2). Nous nous proposons de prouver le théorème suivant, analogue de 1.1 :

\textbf{Théorème 2.1.} Soient \(S\) un schéma quasi-compact, \(E\) un \(\mathcal O_S\)-Module localement libre de rang \(r+1\), \(X=\mathbf P(E)\) le fibré projectif associé ; soit \(\xi\) la classe dans \(K^{\bullet}(X)\) du complexe égal à \(\mathcal O_X(1)\) en degré nul, et nul pour les autres degrés. Alors \(K^{\bullet}(X)\), muni de sa structure naturelle de \(K^{\bullet}(S)\)-algèbre est un \(K^{\bullet}(S)\)-module libre admettant pour base \(1,\xi,\ldots,\xi^r\).

\textbf{2.2.} Fixons les notations : \(B\) désignera l'algèbre symétrique de \(E\); on notera \(\mathcal A\) la catégorie des \(\mathcal O_X\)-Modules, \(\mathcal G\) la catégorie des \(B\)-Modules gradués avec pour morphismes les morphismes gradués, \(\mathcal G_0\) la sous-catégorie pleine de \(\mathcal G\) dont les objets sont les \(B\)-Modules localement libres de type fini en tant que \(B\)-Modules gradués, \(\mathcal G'\) la catégorie des \(\mathcal O_S\)-Modules. Les catégories \(\mathcal A\), \(\mathcal G\), et \(\mathcal G'\) sont abéliennes. Pour tout \(n\), et tout complexe \(L^\bullet\) de \(D^+(\mathcal A)\), nous noterons \(L^\bullet(n)\) le complexe \(L^\bullet\otimes_{\mathcal O_X}\mathcal O_X(n)\).

\textbf{2.3.}
